\section*{Contents}
\begin{tabularx}{0.98\linewidth}{@{}>{\bfseries}p{0.28\linewidth}X@{}}
  1 & Frontier Models --- 最前線の料金表と提供条件 \\
  2 & Open Reasoning --- 公開ウェイトの推論モデル \\
  3 & Decision Inference --- 判断支援の3界面 \\
  4 & DevDay Surface --- 働く窓口 \\
  5 & Safety and Provenance --- 安全性と来歴の5領域 \\
  6 & Training Methods and Systems --- 学習の仕組み \\
  7 & Evaluation Infrastructure --- 評価基盤 \\
  8 & Multimodal Serving --- マルチモーダルの提供現場 \\
  9 & Enterprise Digest --- 企業・業界動向 \\
  10 & Week in Review --- 変化・意味・次の観測点 \\
\end{tabularx}

\vspace{4mm}
\section*{This Week in AI}

\begin{tcolorbox}[enhanced,colback=SurveySoft,colframe=SurveyRule,boxrule=0.4pt,arc=1mm,left=3mm,right=3mm,top=2mm,bottom=2mm]
\textbf{最前線には料金表が並んだ。}\quad 9月28日のSonnet 5.5、9月29日のGPT-6.1 Sol、9月30日のGemini 4 Argonに、公開ウェイトの推論モデル、判断支援の3界面、働く窓口、5領域の安全性、学習の仕組み、評価基盤、マルチモーダルの提供現場、企業・業界動向が重なった。
\end{tcolorbox}

\begin{tcolorbox}[enhanced,colback=SurveySoft,colframe=SurveyRule,boxrule=0.4pt,arc=1mm,left=3mm,right=3mm,top=2mm,bottom=2mm]
\textbf{数値と日付は出どころごとに受け止める。}\quad ベンチマークの値はベンダー計測によるもので、独立した再現はない。料金表は発表時点の公式ページであり、請求額の実測ではない。記事の日付と評価スクリプトの到着、重みやコードの公開は別の事柄として切り分ける。未検証の事柄、査読前の報告、構想段階のものは、その旨を添えて読む。
\end{tcolorbox}

\begin{tcolorbox}[enhanced,colback=SurveySoft,colframe=SurveyRule,boxrule=0.4pt,arc=1mm,left=3mm,right=3mm,top=2mm,bottom=2mm]
\textbf{載せないものは理由とともに記す。}\quad 素材的に重い2件のホスト記事は、一次資料の受け入れ整理が終わっていないため今号の章には載せない。締め切り後の速報は次号へ送り、通常の集計には含めない。公開投稿の反応は関心の方向を示す背景にとどめる。
\end{tcolorbox}

\vspace{2mm}
\begin{claimboundary}[この号の読み方]
通常本文の対象期間は2026-09-25 18:00から2026-10-02 18:00（いずれもAmerica/New\_York、終了時刻は含まない）。この範囲の出来事を通常の扱いとし、期間前の持ち越し、締め切り後の速報、予告された今後の運用は区別する。本号で参照する公開資料が示す範囲を超えて断定しない。ベンダーや報道が伝える数値はそれぞれの主体に帰属させ、条件の違う測定値を横断的な順位へ言い換えない。モデルの公開場所や一次発表で確認できない仕様や利用条件は事実として扱わない。
\end{claimboundary}

\vspace{2mm}
\noindent\textit{読み解く軸：} 料金と提供形態、公開ウェイトとライセンスの切り分け、ベンダー計測の帰属、記事の日付とスクリプト到着の切り分けを、一次資料の範囲と留保とともに束ねることである。
